% ---------------------------------------------------------------------
%  Chapter 5 : Colored Noise and the Markovian Embedding
% ---------------------------------------------------------------------
\chapter{Colored Noise and the Markovian Embedding}\label{ch:colored-markovian}

Every model in the previous two chapters drove its fields with \emph{white}
noise --- a random force that is uncorrelated across any two distinct instants.
White noise is the easy case, and not by accident: the temporal integrators that
do the real work at the end of the pipeline have fast, fully \emph{analytic}
mode-sum routines that exist precisely because the noise two-point function is a
Dirac delta in time. The moment you ask for a noise with \emph{memory} --- a
finite correlation time $\tau_c$, a smooth kernel like $\ee^{-|\tau|/\tau_c}$ ---
those analytic routines no longer apply, and the pipeline is forced onto a
generic numerical quadrature that, at one loop and beyond, either crawls or
hangs.

This chapter is about the one preprocessor that makes colored noise tractable
anyway. It is a \term{builder-level rewrite}: before your model is ever compiled,
it silently replaces your colored-noise term with an equivalent \emph{white}-noise
problem by introducing one extra dynamical field --- an auxiliary
Ornstein--Uhlenbeck (OU) process --- per noisy field. The output statistics are
identical; the representation is larger but every integral in it is again the
fast white-noise kind. The whole subsystem lives in one file,
\file{pipeline/colored\_to\_markovian.py}, and this chapter walks it from the
physics down to the exact text strings it emits.

It sits in Part~II (\emph{Specifying a Model}) rather than in the
diagram-evaluation parts because it happens entirely at authoring time: it is a
transformation \emph{of the model spec}, run inside \code{TheoryBuilder.build()}
before any symbolic compilation. By the time the field-theory core (Chapter~6
onward) sees your model, the colored noise is gone and only white noise remains.

\section{Motivation: why colored noise breaks the fast integrators}

A noisy field $x(t)$ is driven by a random force $\xi(t)$. For a Gaussian noise,
that force is fully characterised by its mean (zero) and its two-time
autocovariance --- the \emph{second cumulant}:
\begin{equation}
  \avg{\xi(t)\,\xi(t')} \;=\; \kappa(t-t') \;=\; \kappa(\tau),
  \qquad \tau = t - t'.
  \label{eq:cm-secondcumulant}
\end{equation}
Two regimes matter here.

\begin{description}
  \item[White noise.] $\kappa(\tau) = 2D\,\delta(\tau)$. The force is
        uncorrelated across any two distinct instants. This is what the fast
        analytic integrators handle: a delta in time collapses one of the loop
        integrals exactly, leaving a closed-form mode sum.
  \item[Colored noise (single Lorentzian).] $\kappa(\tau) = c\,\ee^{-|\tau|/\tau_c}$.
        The force ``remembers'' itself for a correlation time $\tau_c$; as
        $\tau_c \to 0$ with $c/\tau_c$ held fixed it limits back to white. The
        name \term{Lorentzian} is because the Fourier transform of
        $\ee^{-|\tau|/\tau_c}$ is a Lorentzian, $\propto 1/(\omega^2 + 1/\tau_c^2)$.
\end{description}

The trouble is the smooth factor. \Daedalus's temporal integrator ---
\code{msrjd.integration.time\_domain.final\_integral} --- cannot absorb a smooth
$\ee^{-|\tau|/\tau_c}$ into its analytic mode-sum routines. The module docstring
states the pain exactly (\file{pipeline/colored\_to\_markovian.py:6--10}):

\begin{quote}\itshape
\code{msrjd.integration.time\_domain.final\_integral} cannot absorb a smooth
\code{exp(-|tau|/tauc)} factor coming from a CGF-tab \code{NoiseSourceType}
vertex into its analytic mode-sum integrators. At \code{max\_ell >= 1} the scipy
\code{nquad} fallback path either hangs or runs for orders of magnitude longer
than the equivalent white-noise theory.
\end{quote}

Here \code{nquad} is SciPy's nested numerical multi-dimensional quadrature --- the
slow generic fallback. And \code{max\_ell} is the loop-order cutoff $\ell$:
$\ell = 0$ is tree level, $\ell \ge 1$ adds the loop corrections. So the bottom
line is blunt: at $\ell \ge 1$ a colored theory is effectively unusable.

\begin{note}
This is purely a \emph{performance} cliff, not a correctness one. At tree level
($\ell = 0$) a colored theory will still run through the \code{nquad} fallback
and give the right answer. The embedding in this chapter exists so that the
\emph{loop} corrections become affordable --- it turns an intractable
quadrature into the same fast analytic mode sum every white-noise model uses.
\end{note}

\subsection{How a colored cumulant is declared}

In \Daedalus{} a colored second cumulant is declared as a \term{CGF row} ---
a cumulant-generating-functional term --- with \code{order=2}, a
\code{coefficient} text (the amplitude $c$), and a \code{kernel} text (the
$\tau$-shape). The canonical example, the OU process with a cubic restoring force
and colored driving, declares it like this
(\file{theories/ou\_quartic\_colored.theory.py:18}):

\begin{lstlisting}
.declare_cgf_term('CXX', response_legs=['xt', 'xt'], order=2,
                  coefficient='2*D/tauc', kernel='exp(-abs(tau)/tauc)')
\end{lstlisting}

so here $c = 2D/\tau_c$ and the kernel is $\ee^{-|\tau|/\tau_c}$ --- the full
second cumulant is $(2D/\tau_c)\,\ee^{-|\tau|/\tau_c}$. The response legs
\code{['xt', 'xt']} say the cumulant connects the response field \code{xt}
(the \msrjd{} conjugate of \code{x}) to itself: it is an \emph{auto}-cumulant
on a single field. The accompanying action text is the plain deterministic
dynamics,
\begin{lstlisting}
.set_action_text('xt*((Dt+mu)*x + eps*x^3)')
\end{lstlisting}
--- notice there is \emph{no} noise vertex in the action text. For colored noise
the noise lives entirely in the CGF row, not as a $-D\tilde x^2$ term in the
action (contrast the white-noise OU quartic of Chapter~2, whose action carried
\code{- D*xt[i]\^{}2} explicitly).

\begin{defn}
A \term{CGF row} is a declared contribution to the noise's cumulants, given as a
4-tuple \code{(coefficient, kernel, order, response\_legs)}. \code{order=2} is the
noise \emph{covariance}. The \code{coefficient} is the amplitude, the
\code{kernel} is the function of the lag $\tau$; \code{response\_legs} names which
response fields the cumulant connects. Colored noise sets a smooth kernel like
\code{exp(-abs(tau)/tauc)}; white noise sets \code{dirac\_delta(tau)}.
\end{defn}

\section{The fix in one idea: a colored noise is a filtered white noise}

There is a classic trick in stochastic dynamics. A colored noise whose
autocovariance is a single decaying exponential is \emph{exactly} the output of a
first-order linear filter --- an Ornstein--Uhlenbeck process --- driven by white
noise. So instead of feeding colored noise $\xi(t)$ directly into your equation,
you do two things:

\begin{enumerate}
  \item drive $\xi$ \emph{itself} with white noise through an OU relaxation
        equation, and
  \item feed the OU \emph{output} $\xi$ into your field as an ordinary linear
        coupling.
\end{enumerate}

Your field now sees a deterministic linear coupling to a new auxiliary field, and
the only stochastic forcing left anywhere in the system is white. The white-noise
analytic integrators apply again. The price is one extra field (and its \msrjd{}
response partner) per noisy field: a bigger theory, but a cheaper-to-integrate
one. The module docstring writes the two-line system out
(\file{pipeline/colored\_to\_markovian.py:13--23}):
\begin{align}
  \frac{\dd x}{\dd t}  \;&=\; f(x) \;+\; \xi,
  \label{eq:cm-toptline}\\[2pt]
  \frac{\dd \xi}{\dd t} \;&=\; -\frac{\xi}{\tau_c}
        \;+\; \sqrt{\tfrac{2c}{\tau_c}}\;\eta(t),
        \qquad \avg{\eta(t)\,\eta(t')} = \delta(t-t').
  \label{eq:cm-bottomline}
\end{align}

Read this as a coupled pair:

\begin{description}
  \item[Top line, \eqref{eq:cm-toptline}] is \emph{your} equation, except the
        colored force $\xi$ is no longer ``random by fiat'' --- it is now a
        genuine dynamical field.
  \item[Bottom line, \eqref{eq:cm-bottomline}] is the \emph{auxiliary} OU
        equation: $\xi$ relaxes toward zero with rate $1/\tau_c$, and is itself
        kicked by \emph{white} noise $\eta$ of unit strength, scaled by the
        amplitude $\sqrt{2c/\tau_c}$.
\end{description}

\begin{defn}
The \term{Ornstein--Uhlenbeck (OU) process} is the simplest mean-reverting linear
SDE, $\dd\xi/\dd t = -\xi/\tau_c + (\text{white})$. The \term{Markovian
embedding} is the act of enlarging the state space (adding auxiliary fields) so
that a non-Markovian, memory-carrying, colored process becomes Markovian
(memoryless, white-driven). The observable output field is unchanged; only the
representation grows.
\end{defn}

\subsection{Why the OU output is exactly the Lorentzian we wanted}

The reason the whole construction works is a single claim: the \emph{stationary}
autocovariance of the OU output $\xi$ defined by \eqref{eq:cm-bottomline} is
exactly the Lorentzian we set out to reproduce,
\begin{equation}
  \avg{\xi(t)\,\xi(t')} \;=\; c\,\ee^{-|t-t'|/\tau_c}.
  \label{eq:cm-ouvariance}
\end{equation}

\emph{Derivation.} The OU SDE $\dd\xi/\dd t = -\xi/\tau_c + \sigma\eta$, with
white $\eta$ ($\avg{\eta(t)\eta(t')} = \delta(t-t')$) and
$\sigma = \sqrt{2c/\tau_c}$, is linear, so its stationary two-point function is
known in closed form. The stationary \emph{variance} balances the injection rate
$\sigma^2$ against the dissipation rate $2/\tau_c$, giving $\sigma^2 \tau_c/2$;
the decay in the lag is governed by the relaxation rate $1/\tau_c$. Hence
\begin{equation}
  \avg{\xi\,\xi}(\tau)
    \;=\; \frac{\sigma^2 \tau_c}{2}\,\ee^{-|\tau|/\tau_c}
    \;=\; \frac{(2c/\tau_c)\,\tau_c}{2}\,\ee^{-|\tau|/\tau_c}
    \;=\; c\,\ee^{-|\tau|/\tau_c}.
  \label{eq:cm-ouvariance-derived}
\end{equation}
That last equality is exactly \eqref{eq:cm-ouvariance}. \hfill$\checkmark$

The derivation also pins the white drive amplitude. In \emph{variance} terms, the
strength of the white kick is $\sigma^2 = 2c/\tau_c$; \Daedalus{} declares white
noise as a CGF row with kernel $\delta(\tau)$, and the coefficient of such a row
is exactly the variance strength. So the auxiliary white-noise coefficient must be
\begin{equation}
  \text{aux white-noise coefficient} \;=\; \frac{2c}{\tau_c}.
  \label{eq:cm-auxdrive}
\end{equation}
This $2c/\tau_c$ is precisely what the detector computes and stores as the
\code{aux\_drive\_text}. The code builds it symbolically
(\file{pipeline/colored\_to\_markovian.py:183}):

\begin{lstlisting}
aux_drive_expr = (SR(2) * c_expr / tauc_expr).simplify_full()
\end{lstlisting}

For the worked example $c = 2D/\tau_c$, this gives
$2\cdot(2D/\tau_c)/\tau_c = 4D/\tau_c^2$, which the test suite pins down
(\file{tests/test\_markovianize.py:73}):
\code{assert info['aux\_drive\_text'] == '4*D/tauc\^{}2'}.

\section{From SDE to MSR-JD action: the three structural pieces}

\Daedalus{} works in the \msrjd{} (Martin--Siggia--Rose / Janssen--De~Dominicis)
path integral. Each physical field $x$ gets a conjugate \term{response field}
$\tilde x$ (written \code{xt} in code), and a linear equation of motion
$(\Dt + \mu)x = (\text{forcing})$ appears in the action as the bilinear term
$\tilde x\,[(\Dt + \mu)x - (\text{forcing})]$. The preprocessor's edits are
nothing more than the \msrjd{} encoding of the two SDE lines
\eqref{eq:cm-toptline}--\eqref{eq:cm-bottomline}. There are exactly three
structural pieces, plus two text edits to the action.

\paragraph{Piece 1 --- the auxiliary field.} For each noisy field, inject one new
physical field $\xi$ (code name \code{xi}). \Daedalus's authoring API auto-creates
its conjugate response \code{xit} and its saddle parameter \code{xistar}, just as
for any user-declared field.

\paragraph{Piece 2 --- the auxiliary OU drift equation.} The OU line
\eqref{eq:cm-bottomline}, stripped of its white drive, is the deterministic
operator $(\Dt + 1/\tau_c)\xi = 0$. The code injects this equation directly. In
the action it becomes a kinetic term contracted with the aux response field
\code{xit} (\file{pipeline/colored\_to\_markovian.py:431--435}):
\begin{lstlisting}[style=console]
    + xit*((Dt + 1/tauc)*xi)
\end{lstlisting}

\paragraph{Piece 3 --- the auxiliary white CGF row.} The white drive
$\sqrt{2c/\tau_c}\,\eta$ does \emph{not} appear in the action text. It is
represented by a \emph{new white CGF row} on the aux response legs --- the
$2c/\tau_c$ coefficient of \eqref{eq:cm-auxdrive} with kernel $\delta(\tau)$ ---
which is how \Daedalus{} injects Gaussian white noise everywhere. This is the
\code{new\_cgf\_rows} construction at
\file{pipeline/colored\_to\_markovian.py:397--405}.

That leaves one text edit to wire the auxiliary output into your dynamics:

\paragraph{Edit --- couple the user's field to the aux output.} The colored
force $\xi$ enters the SDE as $\dd x/\dd t = f(x) + \xi$, i.e. moving everything to
one side, $(\Dt + \mu)x + \varepsilon x^3 - \xi = 0$. So the code inserts
\code{- xi} inside your parenthesised dynamics block
(\file{pipeline/colored\_to\_markovian.py:472}):
\begin{lstlisting}[style=console]
    xt*((Dt+mu)*x + eps*x^3)   ->   xt*((Dt+mu)*x + eps*x^3 - xi)
\end{lstlisting}

Putting the three pieces together, the auxiliary field, the aux OU drift
equation, and the aux white CGF row jointly reproduce the bottom SDE line
$\dd\xi/\dd t = -\xi/\tau_c + \sqrt{2c/\tau_c}\,\eta$; the \code{- xi} coupling
plus the appended $\tilde\xi\,(\Dt + 1/\tau_c)\,\xi$ kinetic term wire $\xi$ into
the top line. The model that emerges is a pure white-noise theory.

\begin{note}
Note the asymmetry between the two text edits. The \code{- xi} \emph{coupling} is
inserted \emph{inside} an existing parenthesised block (it must land next to the
right field's dynamics), whereas the aux \emph{kinetic} term is simply
\emph{appended} to the end of the action string with a leading \code{+}. Appending
is safe because the kinetic term is self-contained; the coupling insertion is the
delicate one, and \S\ref{sec:cm-injectcoupling} is where its fragility lives.
\end{note}

\section{The cross-correlated multi-field case}

The embedding generalises cleanly to several fields whose colored forces are
\emph{cross}-correlated. For two fields $x, y$ with
\begin{align}
  \avg{\xi_x\,\xi_x}(\tau) &= c_{xx}\,\ee^{-|\tau|/\tau_c}, &
  \avg{\xi_y\,\xi_y}(\tau) &= c_{yy}\,\ee^{-|\tau|/\tau_c}, &
  \avg{\xi_x\,\xi_y}(\tau) &= c_{xy}\,\ee^{-|\tau|/\tau_c},
  \label{eq:cm-cross}
\end{align}
the embedding gives $x$ its own auxiliary OU process $\xi_x$ (code \code{xi}) and
$y$ its own $\xi_y$ (code \code{yi}), each driven by white noise, plus a single
\emph{cross} white CGF row between the two aux response legs. The shared
relaxation rate $1/\tau_c$ is what makes this work: because all three covariances
in \eqref{eq:cm-cross} decay at the same rate, the aux white covariance matrix
between $(\eta_x, \eta_y)$ factorises cleanly and reproduces the desired colored
cross-covariance.

The shipped two-field example
(\file{theories/ou\_quartic\_two\_dim\_color\_corr.theory.py:25--27}) declares
exactly this:
\begin{lstlisting}
.declare_cgf_term('Cxx', response_legs=['xt','xt'], order=2,
                  coefficient='D1/tauc',           kernel='exp(-abs(tau)/tauc)')
.declare_cgf_term('Cyy', response_legs=['yt','yt'], order=2,
                  coefficient='D2/tauc',           kernel='exp(-abs(tau)/tauc)')
.declare_cgf_term('Cxy', response_legs=['xt','yt'], order=2,
                  coefficient='rho*sqrt(D1*D2)/tauc', kernel='exp(-abs(tau)/tauc)')
\end{lstlisting}
The three colored rows \code{Cxx}, \code{Cyy}, \code{Cxy} become three white
aux rows: two auto rows on \code{[xit,xit]} and \code{[yit,yit]}, plus one cross
row on \code{[xit,yit]}. The docstring lays this out at
\file{pipeline/colored\_to\_markovian.py:33--39}.

\section{The kernel matcher: \texttt{detect\_lorentzian}}\label{sec:cm-detect}

Everything above is the physics. The rest of the chapter is the code that
implements it, starting with the function that decides whether a given CGF row
\emph{is} a single Lorentzian and, if so, extracts $\tau_c$, the amplitude $c$,
and the aux drive $2c/\tau_c$. That function is \code{detect\_lorentzian}
(\file{pipeline/colored\_to\_markovian.py:79}):

\begin{lstlisting}
def detect_lorentzian(kernel_text: str,
                      coefficient_text: str,
                      declared_params: list[str]) -> Optional[dict]:
\end{lstlisting}

It takes the row's \code{kernel} string (e.g. \code{'exp(-abs(tau)/tauc)'}), its
\code{coefficient} string (e.g. \code{'2*D/tauc'}), and the list of declared
parameter names (so it can build a namespace in which \code{tauc}, \code{D}, etc.
are known symbols). It returns a dict describing the embedding on a match, or
\code{None} if the kernel does not cleanly match the single-Lorentzian template.

This is the one part of the subsystem that does real symbolic algebra, so before
reading the steps we pause on the single library it leans on.

\begin{defn}
\term{SageMath} is a large open-source mathematics system built on top of Python;
it bundles many math libraries behind one Python-friendly API. The piece this
module uses is Sage's \term{Symbolic Ring}, abbreviated \code{SR} --- a
``calculator'' that works with algebraic expressions containing \emph{named
variables} (like \code{tauc}, \code{D}, \code{tau}) rather than just numbers. An
element of \code{SR} is a symbolic expression, e.g. \code{exp(-abs(tau)/tauc)},
that you can differentiate, simplify, substitute into, and pretty-print
symbolically. This is the same \emph{idea} as \code{sympy} (which \Daedalus{}
uses elsewhere), but Sage's \code{SR} is a separate, more powerful engine; the two
are \emph{not} interchangeable objects.
\end{defn}

The import line is (\file{pipeline/colored\_to\_markovian.py:70--73}):

\begin{lstlisting}
from sage.all import (
    SR, sage_eval, exp, abs_symbolic, simplify,
    sqrt as sage_sqrt,
)
\end{lstlisting}

What each name is, and exactly how this file uses it:

\begin{description}
  \item[\code{SR}] the Symbolic Ring object. \code{SR.var('name')} creates (and
        globally registers) a symbolic variable; \code{SR(x)} \emph{coerces} an
        arbitrary object $x$ into a symbolic expression so that symbolic methods
        can be called on it. Both forms appear throughout the file.
  \item[\code{sage\_eval}] Sage's \emph{safe, namespaced} string-to-expression
        evaluator. You hand it a string of Sage syntax plus a dictionary mapping
        names to Sage objects, and it parses and evaluates the string \emph{in
        that namespace only}. It is like Python's built-in \code{eval()} but tuned
        to Sage syntax (\code{\^{}} means power) and scoped to the dict you pass.
  \item[\code{exp}] Sage's symbolic exponential. It plays two roles: it is put
        \emph{into} the evaluation namespace so the kernel string's \code{exp(...)}
        resolves to it, and it is \emph{compared against} to recognise the kernel's
        top-level operation.
  \item[\code{abs\_symbolic}] Sage's symbolic absolute value. Plain Python
        \code{abs} would try to evaluate numerically; \code{abs\_symbolic} keeps
        $|\tau|$ as an \emph{unevaluated} symbolic object so the matcher can reason
        about it.
  \item[\code{simplify}, \code{sqrt as sage\_sqrt}] imported but \emph{not} called
        directly in this file. The code uses the more aggressive \emph{method}
        \code{.simplify\_full()} instead of the free \code{simplify}, and the
        \code{sqrt} appearing in coefficient strings is resolved by
        \code{sage\_eval} against Sage's default namespace rather than via this
        alias.
\end{description}

The matcher also calls a handful of methods on \code{SR} expressions:
\code{.operator()} returns the outermost operation node (for a Lorentzian that
must be \code{exp}); \code{.operands()} returns the arguments of that operation
($\exp(z)$ has exactly one operand, $z$); \code{.simplify\_full()} returns a
maximally-simplified copy; and \code{.has(sub)} is a boolean test for whether an
expression contains a given sub-expression anywhere in its tree.

\subsection{Step by step}

\begin{enumerate}
  \item \textbf{Empty guard} (\file{pipeline/colored\_to\_markovian.py:110}). If
        \code{kernel\_text} is falsy, return \code{None}.

  \item \textbf{Build the evaluation namespace}
        (\file{pipeline/colored\_to\_markovian.py:128--135}). Create a
        \emph{fresh local} symbol \code{tau\_sym = SR.var('\_lorentz\_tau')} but
        bind it to the \emph{name} \code{'tau'} in the namespace, so the kernel
        string's \code{tau} resolves to it. Add \code{exp} and \code{abs ->
        abs\_symbolic}, then add every declared parameter as a generic
        \code{SR.var(pname)} via \code{setdefault} (so a parameter the user
        already declared wins, and any other name is still defined):
\begin{lstlisting}
tau_sym = SR.var('_lorentz_tau')
eval_ns: dict = {
    'tau': tau_sym,
    'exp': exp,
    'abs': abs_symbolic,
}
for pname in declared_params:
    eval_ns.setdefault(pname, SR.var(pname))
\end{lstlisting}

  \item \textbf{Parse the kernel}
        (\file{pipeline/colored\_to\_markovian.py:137--141}).
        \code{K = sage\_eval(kernel\_text, eval\_ns)} inside a \code{try/except}
        that returns \code{None} on any parse error, then \code{K = SR(K)} to
        guarantee a symbolic object.

  \item \textbf{Top-level must be \texttt{exp}}
        (\file{pipeline/colored\_to\_markovian.py:143--148}). If
        \code{K.operator() != exp}, return \code{None}. Then read
        \code{operands = K.operands()}, require exactly one, take
        \code{arg = SR(operands[0])} (the exponent), and build
        \code{abs\_tau = abs\_symbolic(tau\_sym)}.

  \item \textbf{Decompose} $\mathtt{arg} = -|\tau|/\tau_c$
        (\file{pipeline/colored\_to\_markovian.py:154--161}). Compute
        \code{r = (arg / abs\_tau).simplify\_full()}. For a true Lorentzian, $r$
        should be $-1/\tau_c$ --- which is to say \emph{abs-free} and
        \emph{tau-free}. The code rejects (\code{return None}) if
        \code{r.has(abs\_tau)} or \code{r.has(tau\_sym)}. This single test is what
        distinguishes the three cases below.

  \item \textbf{Compute and validate} $\tau_c$
        (\file{pipeline/colored\_to\_markovian.py:163--170}).
        \code{tauc\_expr = (SR(-1) / r).simplify\_full()}, then reject if it still
        references \code{tau\_sym} or \code{abs\_tau}.

  \item \textbf{Parse the amplitude}
        (\file{pipeline/colored\_to\_markovian.py:174--180}).
        \code{c\_expr = sage\_eval(coefficient\_text or '0', eval\_ns)}, coerced
        and \code{.simplify\_full()}-ed; \code{return None} on failure. The
        \code{or '0'} guards an empty coefficient string. This is the $c$ in
        $c\,\ee^{-|\tau|/\tau_c}$.

  \item \textbf{Compute the aux drive}
        (\file{pipeline/colored\_to\_markovian.py:183}). The
        \code{aux\_drive\_expr = (SR(2) * c\_expr / tauc\_expr).simplify\_full()}
        from \eqref{eq:cm-auxdrive}.

  \item \textbf{Return the embedding dict}
        (\file{pipeline/colored\_to\_markovian.py:185--192}) with both the SR
        expressions and their string renderings (via \code{\_sr\_to\_text},
        \S\ref{sec:cm-srtotext}): \code{tauc\_expr}, \code{tauc\_text},
        \code{amplitude\_expr}, \code{amplitude\_text}, \code{aux\_drive\_expr},
        \code{aux\_drive\_text}.
\end{enumerate}

The decomposition in step~5 is the whole matcher in miniature. Three kernels and
what $r = \mathtt{arg}/|\tau|$ becomes for each:

\begin{center}
\begin{tabular}{@{}lll@{}}
\toprule
kernel & $r = \mathtt{arg}/|\tau|$ & verdict\\
\midrule
$\ee^{-|\tau|/\tau_c}$  & $-1/\tau_c$ (abs-free, tau-free)        & \textbf{match} \\
$\ee^{-\tau/\tau_c}$    & $-\tau/(\tau_c\,|\tau|)$ (still has $\tau$) & reject \\
$\ee^{-\tau^2/(\cdots)}$ & still has $\tau$                       & reject \\
\bottomrule
\end{tabular}
\end{center}

The middle row is the physically important rejection: a bare exponential
$\ee^{-\tau/\tau_c}$ without the absolute value is \emph{asymmetric} in $\tau$,
and a true autocovariance must be symmetric (it depends on $\tau = t-t'$ only
through $|\tau|$). The matcher correctly refuses it. The test suite checks all
three verdicts directly, including the Gaussian and the bare-exponential
rejections (\file{tests/test\_markovianize.py}, the
\code{test\_lorentzian\_kernel\_detection} case).

\begin{gotcha}
\textbf{The \texttt{tau} pollution guard.} Step~2 deliberately mints a fresh
\code{SR.var('\_lorentz\_tau')} instead of reusing the global \code{tau}. The long
comment at \file{pipeline/colored\_to\_markovian.py:116--127} explains why. Other
parts of the pipeline routinely register a global \code{tau} symbol carrying a
\code{domain='positive'} assumption (for instance an exponential-synapse time
constant). If $\tau > 0$ is assumed globally, then \code{abs(tau)} collapses to
\code{tau}, and the matcher would \emph{wrongly accept} an asymmetric kernel like
$\ee^{-\tau/\tau_c}$ --- precisely the case it is supposed to reject. Using a name
the caller never declares positive keeps the matcher robust against
test-ordering pollution. This is subtle and order-dependent: a regression would
only surface if a positive \code{tau} had been registered earlier in the same
Sage session.
\end{gotcha}

\begin{gotcha}
\textbf{Positivity of $\tau_c$ is documented but not enforced.} The docstring
says $\tau_c$ ``must be POSITIVE for a physical Lorentzian'' and that the matcher
requires the decay rate to ``evaluate to a strictly positive expression''
(\file{pipeline/colored\_to\_markovian.py:91--93, 163--164}), but the code only
ever checks that $\tau_c$ is abs-free and tau-free --- never that it is
numerically positive. A symbolic $\tau_c$ that happens to be negative for some
parameter values would still be accepted. Keep your $\tau_c$ a genuinely positive
parameter.
\end{gotcha}

\section{The builder transform: \texttt{markovianize\_spec}}\label{sec:cm-transform}

The detector is called from one place: \code{markovianize\_spec}
(\file{pipeline/colored\_to\_markovian.py:198}), the entry point that actually
rewrites the builder. Its signature is deceptively simple:

\begin{lstlisting}
def markovianize_spec(builder) -> None:
\end{lstlisting}

It takes a \code{TheoryBuilder} (in practice a \code{TemporalTheoryBuilder}) and
returns \emph{nothing}. All of its effects are in-place mutations: it reads
\code{builder.\_cgf\_terms}, \code{builder.\_markovianize\_default},
\code{builder.parameters}, \code{builder.physical\_fields}, and
\code{builder.\_action\_text}, and it writes them back --- crucially, by calling
the very same public authoring methods a user would call
(\code{builder.physical\_field(...)}, \code{builder.equation(...)}) plus direct
overwrites of \code{\_cgf\_terms} and \code{\_action\_text}.

\begin{note}
\textbf{Idempotency.} Re-running \code{markovianize\_spec} on an
already-markovianized builder is a no-op: no remaining row matches, because the
injected aux rows carry \code{markovianize=False} and the original colored rows
are gone. This safety property rests entirely on that one flag
(\file{pipeline/colored\_to\_markovian.py:404}); if it were dropped, a second
\code{build()} pass could try to re-embed the white aux rows.
\end{note}

\subsection{Step by step}

\begin{enumerate}
  \item \textbf{Snapshot \& empty guard}
        (\file{pipeline/colored\_to\_markovian.py:223--225}). Copy the rows;
        return immediately if there are none.

  \item \textbf{Builder-level gating}
        (\file{pipeline/colored\_to\_markovian.py:227--230}). Read
        \code{builder\_default = \_markovianize\_default} (default \code{True}).
        Compute \code{any\_explicit\_on = any(t.get('markovianize') is True ...)}.
        If the default is OFF \emph{and} no row explicitly opts in, return. This
        is the \code{.markovianize(False)} opt-out path.

  \item \textbf{Collect declared parameter names}
        (\file{pipeline/colored\_to\_markovian.py:232}) for the matcher.

  \item \textbf{Classify each row}
        (\file{pipeline/colored\_to\_markovian.py:240--276}) into
        \code{matched\_rows} / \code{passthrough\_rows}:
        \begin{itemize}
          \item \code{markovianize is False} $\to$ passthrough.
          \item \code{order != 2} $\to$ passthrough (the colored embedding is a
                \emph{second}-cumulant story).
          \item otherwise run \code{detect\_lorentzian(...)}.
          \item on \textbf{no match}: if the row had \code{markovianize is True}
                (an \emph{explicit} assertion), \textbf{raise \code{ValueError}}
                telling the user their kernel doesn't match the template;
                otherwise pass through silently.
          \item on \textbf{match}: append \code{\{'term': term, 'info': info\}} to
                \code{matched\_rows}.
        \end{itemize}

  \item \textbf{Early out}
        (\file{pipeline/colored\_to\_markovian.py:278--279}). If nothing matched,
        return.

  \item \textbf{Gather response-leg fields}
        (\file{pipeline/colored\_to\_markovian.py:292--298}). Read each matched
        row's \code{response\_legs}; if absent, broadcast its
        \code{response\_field} to a 2-list; collect all leg names into
        \code{source\_field\_names}.

  \item \textbf{Invert the response$\to$natural map}
        (\file{pipeline/colored\_to\_markovian.py:303--317}). For each response
        field name \code{rname} (e.g. \code{'xt'}), find the physical field whose
        auto-response equals it by testing \code{f'\{nat\}t' == rname} (recall
        \code{physical\_field('x')} auto-creates the response \code{xt}); store
        \code{response\_to\_natural[rname] = nat} (e.g. \code{'x'}). If none is
        found, store \code{None} --- that row will fall back to scipy because it
        cannot be embedded.

  \item \textbf{Single-$\tau_c$ enforcement}
        (\file{pipeline/colored\_to\_markovian.py:323--334}). Collect the set of
        distinct \code{tauc\_text} strings across matched rows. If there is more
        than one, \textbf{raise \code{NotImplementedError}} --- v1 supports a
        single shared $\tau_c$. Else take the one $\tau_c$.

  \item \textbf{Assign aux field names}
        (\file{pipeline/colored\_to\_markovian.py:339--352}). For each source
        natural name, in sorted order for determinism, call
        \code{\_pick\_aux\_field\_name(...)} (\S\ref{sec:cm-pickname}) to get a
        collision-safe aux name (\code{x -> xi}), record it, and \emph{reserve} it
        against later collisions in the same pass. If the map ends up empty,
        return.

  \item \textbf{Inject aux physical fields}
        (\file{pipeline/colored\_to\_markovian.py:355--360}). For each
        \code{(source\_nat, aux\_nat)}, call
        \code{builder.physical\_field(aux\_nat, indexed=True, description=...)}.
        This is the standard authoring method --- it auto-creates the conjugate
        response \code{<aux\_nat>t} and the saddle parameter
        \code{<aux\_nat>star}.

  \item \textbf{Inject aux OU drift equations}
        (\file{pipeline/colored\_to\_markovian.py:363--367}). For each aux field,
        \code{builder.equation(lhs=f'(Dt + 1/\{tauc\_text\}) * \{aux\_nat\}',
        rhs='0')} --- the deterministic part of the OU process,
        $(\Dt + 1/\tau_c)\xi = 0$.

  \item \textbf{Build replacement white CGF rows}
        (\file{pipeline/colored\_to\_markovian.py:370--405}). For each matched
        row, recover its two response legs, map them to natural names
        \code{nat\_a}, \code{nat\_b}, derive the aux response names
        \code{aux\_resp\_a}, \code{aux\_resp\_b}, and emit one white row:
\begin{lstlisting}
{
    'name':           f'{term["name"]}_markov_aux',
    'response_field': None,
    'response_legs':  [aux_resp_a, aux_resp_b],
    'order':          order,
    'coefficient':    info['aux_drive_text'],   # 2c / tauc
    'kernel':         'dirac_delta(tau)',        # white
    'markovianize':   False,                     # don't re-process
}
\end{lstlisting}
        For an auto-cumulant \code{[xt, xt]} both aux legs are \code{xit} (one
        row). For a cross-cumulant \code{[xt, yt]} the legs are \code{[xit, yit]}
        (the cross row); the auto rows \code{Cxx}, \code{Cyy} separately produce
        \code{[xit,xit]} and \code{[yit,yit]}.

  \item \textbf{Swap the CGF list}
        (\file{pipeline/colored\_to\_markovian.py:408}).
        \code{builder.\_cgf\_terms = passthrough\_rows + new\_cgf\_rows} --- the
        colored rows are gone, replaced by white aux rows; untouched rows are
        preserved.

  \item \textbf{Augment the action text}
        (\file{pipeline/colored\_to\_markovian.py:424--436}). For each source
        field, call \code{\_inject\_aux\_coupling(action\_text, resp\_field,
        aux\_nat)} with \code{resp\_field = f'\{source\_nat\}t'} to insert
        \code{- xi} inside the user's \code{xt*(...)} block. Then for each aux
        field, append \code{+ \{aux\_resp\}*((Dt + 1/\{tauc\_text\})*\{aux\_nat\})}
        --- the OU kinetic term. Write back to \code{builder.\_action\_text}.

  \item \textbf{Record diagnostics}
        (\file{pipeline/colored\_to\_markovian.py:439--443}). Set
        \code{builder.\_markovianize\_applied} to a small dict recording the aux
        field names, the matched row names, and the shared \code{tauc\_text}.
\end{enumerate}

\subsection{Where it is wired into \texttt{build()}}

The transform is invoked from exactly one place,
\file{pipeline/theory.py:1735--1737}, inside \code{build()}:

\begin{lstlisting}
if self._cgf_terms:
    from pipeline.colored_to_markovian import markovianize_spec
    markovianize_spec(self)
\end{lstlisting}

This call is deliberately placed \emph{before} \code{self.\_compile\_text\_declarations()}
(\file{pipeline/theory.py:1742}). Ordering is the whole point: by the time the
downstream compiler turns your text into lambdas, it only ever sees the augmented,
white-noise spec. The colored row never reaches the compiler.

\begin{note}
The two gating knobs you can reach as a user are: \code{.markovianize(False)} on
the builder (turns the whole transform OFF), and \code{markovianize=False} (or
\code{=True}) as a keyword on an individual \code{declare\_cgf\_term} call.
\code{False} skips a row; \code{True} asserts ``this IS a Lorentzian --- fail if
you disagree.'' The default (\code{None}/\code{'auto'}) means ``match if you can,
otherwise pass through silently.'' The opt-out path is exercised directly by the
test suite, which checks that \code{.markovianize(False)} leaves the original
colored row in place and injects no auxiliary field
(\file{tests/test\_markovianize.py}, the
\code{test\_markovianize\_opt\_out\_keeps\_colored\_row} case).
\end{note}

\section{The two helpers}

\subsection{\texttt{\_pick\_aux\_field\_name} --- collision-safe naming}
\label{sec:cm-pickname}

\code{\_pick\_aux\_field\_name(source\_nat, taken, additional\_taken=())}
(\file{pipeline/colored\_to\_markovian.py:449}) picks a name for the auxiliary
field that will not clash with anything already in the model. It unions
\code{taken} with \code{additional\_taken}, then tries the natural candidate
\code{f'\{source\_nat\}i'} (so \code{x -> xi}, \code{y -> yi}); if that is taken it
tries \code{f'xi\_\{source\_nat\}'}; and failing that it loops
\code{f'xi\_\{source\_nat\}\_\{n\}'} for $n = 2, 3, \dots$, which always terminates
because the integer suffix grows without bound.

\subsection{\texttt{\_inject\_aux\_coupling} --- action-text surgery}
\label{sec:cm-injectcoupling}

\code{\_inject\_aux\_coupling(action\_text, resp\_field, aux\_natural)}
(\file{pipeline/colored\_to\_markovian.py:472}) is the one piece of literal
string editing. It inserts \code{- <aux>} inside the parenthesised block that
follows \code{<resp\_field>*(} in the action text, or returns the text unchanged
if that pattern is not found. It is the only place the standard-library
\code{re} (regular expressions) module is used.

\begin{defn}
\code{re} is Python's built-in regular-expression module: it matches text
patterns. A \term{regular expression} is a small pattern language --- \code{\textbackslash{}b}
means ``word boundary,'' \code{\textbackslash{}s*} means ``zero or more whitespace characters,''
and \code{re.escape(s)} quotes any special characters in \code{s} so they match
literally.
\end{defn}

The function works in three moves
(\file{pipeline/colored\_to\_markovian.py:486--508}):
\begin{enumerate}
  \item Compile the pattern (word boundary, the field token, optional whitespace,
        \code{*}, optional whitespace, \code{(}), with \code{re.escape} quoting the
        field name:
\begin{lstlisting}
import re
pattern = re.compile(
    rf'\b{re.escape(resp_field)}\s*\*\s*\(', flags=0)
m = pattern.search(action_text)
\end{lstlisting}
  \item \code{.search(...)} for the first occurrence; if there is none, return the
        text unchanged.
  \item From the match end, \code{start = m.end() - 1} is the index of the
        \code{(}. Walk forward with a \code{depth} counter --- \code{+1} on
        \code{(}, \code{-1} on \code{)} --- until \code{depth} returns to $0$; that
        is the \emph{matching} close paren. Insert \code{f' - \{aux\_natural\}'}
        just before it.
\end{enumerate}
Regex only \emph{locates} the insertion point; the brace matching is a manual
depth counter, because regular expressions cannot reliably match balanced
parentheses. The function edits only the \emph{first} \code{<resp\_field>*(}
occurrence per call; \code{markovianize\_spec} calls it once per source field.

\begin{gotcha}
\textbf{Action-text surgery is fragile by design.} \code{\_inject\_aux\_coupling}
only handles the canonical \code{<rt>*(...)} layout. If your action is written in
a more baroque form, the insertion silently no-ops and the \code{- xi} coupling is
simply \emph{missing} --- you get a silently \emph{wrong} theory, not an error.
There is no validation that the coupling actually landed. This is the most
surprising failure mode in the subsystem. The remedy the docstring prescribes
(\file{pipeline/colored\_to\_markovian.py:418--423, 481--485}) is to opt out with
\code{.markovianize(False)} and hand-roll the embedding yourself.
\end{gotcha}

\subsection{\texttt{\_sr\_to\_text} --- symbolic back to re-parseable string}
\label{sec:cm-srtotext}

\code{\_sr\_to\_text(expr)} (\file{pipeline/colored\_to\_markovian.py:511}) is a
one-liner: \code{return str(SR(expr))}. The default Sage string form uses
\code{\^{}} for power and function names like \code{sqrt}, which is exactly the
syntax that \code{pipeline.theory\_compiler.\_CGFKernelCallable} re-evaluates
later (\file{pipeline/colored\_to\_markovian.py:517--519}). So the round-trip
symbolic~$\to$~text~$\to$~symbolic is lossless: every coefficient the detector
computes can be written into a new CGF row and re-parsed downstream without loss.

\section{Worked example: \texttt{ou\_quartic\_colored} before and after}

Concretely, take the shipped \code{ou\_quartic\_colored} theory and watch the
builder state change across \code{markovianize\_spec}.

\paragraph{Before.} Just before the call, the builder holds:
\begin{itemize}
  \item \textbf{physical fields:} \code{x} (internal \code{dx}, response
        \code{xt}).
  \item \textbf{one CGF row:}
        \code{\{'name': 'CXX', 'response\_legs': ['xt','xt'], 'order': 2,
        'coefficient': '2*D/tauc', 'kernel': 'exp(-abs(tau)/tauc)',
        'markovianize': None\}}.
  \item \textbf{action text:} \code{'xt*((Dt+mu)*x + eps*x\^{}3)'}.
  \item \code{\_markovianize\_default}: \code{True}.
\end{itemize}

\paragraph{Detection.} \code{detect\_lorentzian('exp(-abs(tau)/tauc)', '2*D/tauc',
[...])} returns \code{\{'tauc\_text': 'tauc', 'amplitude\_text': '2*D/tauc',
'aux\_drive\_text': '4*D/tauc\^{}2', ...\}}. The test suite pins both
\code{tauc\_text == 'tauc'} and \code{aux\_drive\_text == '4*D/tauc\^{}2'}
(\file{tests/test\_markovianize.py:71--73}).

\paragraph{After.} The builder now holds:
\begin{itemize}
  \item \textbf{physical fields:} \code{['dx', 'dxi']} --- the auxiliary \code{xi}
        (internal \code{dxi}, auto-response \code{xit}, auto-saddle \code{xistar})
        was added. The test asserts \code{phys\_names == ['dx', 'dxi']}
        (\file{tests/test\_markovianize.py:116}).
  \item \textbf{equations:} the user's $(\Dt+\mu)x = -\varepsilon x^3$ \emph{plus}
        the new aux drift $(\Dt + 1/\tau_c)\xi = 0$ --- two DAE equations.
  \item \textbf{one CGF row}, the colored \code{CXX} replaced by
        \code{\{'name': 'CXX\_markov\_aux', 'response\_legs': ['xit','xit'],
        'order': 2, 'coefficient': '4*D/tauc\^{}2', 'kernel': 'dirac\_delta(tau)',
        'markovianize': False\}}. The test confirms the row name ends in
        \code{\_markov\_aux} and the resolved legs are \code{['xit','xit']}
        (\file{tests/test\_markovianize.py:130, 132}).
  \item \textbf{action text:}
        \code{'xt*((Dt+mu)*x + eps*x\^{}3 - xi) + xit*((Dt + 1/tauc)*xi)'}.
\end{itemize}

\paragraph{Downstream.} \code{make\_correlated\_noises\_block} turns the white aux
row into a \code{model['correlated\_noises']} entry whose order-2 kernel callable
is \code{\_CGFKernelCallable('4*D/tauc\^{}2', 'dirac\_delta(tau)', 2)}. The diagram
engine attaches it to the \code{xit}--\code{xit} noise vertex. Because the kernel
is now $\delta(\tau)$, the fast white-noise analytic integrators run --- which was
the entire goal.

\subsection{The tree-level analytic cross-check}

The payoff is checkable in closed form. With $\varepsilon = 0$ (the free theory),
\code{compute\_cumulants(model, k=2, max\_ell=0, ...)} on the markovianized
\code{ou\_quartic\_colored} must return the analytic colored-OU equal-time
variance
\begin{equation}
  C(0) \;=\; \frac{2D}{\mu\,(\mu\,\tau_c + 1)},
  \label{eq:cm-treevariance}
\end{equation}
the simple-pole result the test pins down
(\file{tests/test\_markovianize.py}, the
\code{test\_colored\_ou\_tree\_level\_matches\_analytic} case, which runs at
$\mu = 0.1,\ \tau_c = 1,\ D = 1$). That the embedded white-noise model reproduces
\eqref{eq:cm-treevariance} is direct confirmation that the OU filter recovers the
original colored statistics.

\section{Limitations and how to live with them}

The v1 embedding is deliberately the minimum that unblocks the working
OU-colored theories. Its boundaries are worth knowing precisely.

\begin{description}
  \item[Single Lorentzian only.] Only $c\,\ee^{-|\tau|/\tau_c}$ matches. Gaussian
        kernels, the bare-exponential $\ee^{-\tau/\tau_c}$ (asymmetric),
        double-exponential, oscillatory, and polynomially-modulated kernels all
        fail the matcher and pass through to the slow scipy fallback. The future
        templates are enumerated in the docstring
        (\file{pipeline/colored\_to\_markovian.py:51--64}).

  \item[Order-2 only.] Only second cumulants are embedded; an \code{order != 2}
        row is passed through silently
        (\file{pipeline/colored\_to\_markovian.py:250--252}). A colored
        \emph{third} cumulant would hit the slow path.

  \item[Single $\tau_c$.] If matched rows carry \emph{different} $\tau_c$
        expressions, v1 raises \code{NotImplementedError}
        (\file{pipeline/colored\_to\_markovian.py:323--333}). All shipped examples
        use one shared $\tau_c$.

  \item[The \texorpdfstring{$\mu = 1/\tau_c$}{mu = 1/tauc} double pole.] This is
        the constraint most likely to bite. See the dedicated box below.
\end{description}

\begin{gotcha}
\textbf{Keep $\mu \neq 1/\tau_c$.} The free ($\ell = 0$) two-point function of the
embedded \emph{user} field is the convolution of the user-field propagator
$1/(\ii\omega + \mu)$ with the OU output spectrum $\propto 1/(\omega^2 +
1/\tau_c^2)$. In the frequency-domain residue calculus this product has poles at
$\omega = \ii\mu$ and $\omega = \ii/\tau_c$, and the analytic mode-sum integrator
computes residues \emph{assuming these poles are simple} (multiplicity~1). When
$\mu = 1/\tau_c$ exactly, the two poles \emph{collide} into a single double pole,
and the simple-pole residue infrastructure cannot form the multiplicity-2
residue. The closed-form variance \eqref{eq:cm-treevariance} is itself the
simple-pole result; at $\mu = 1/\tau_c$ it is the limit of a $0/0$. This is a
\emph{downstream} limitation --- it lives in the residue calculus, not in this
module --- and the embedding code does \emph{not} guard against it. The test
deliberately picks $\mu = 0.1,\ \tau_c = 1$ to avoid the degeneracy, with an
explicit inline note (\file{tests/test\_markovianize.py:165--168}). When you
parameterise a markovianized theory, keep $\mu$ away from $1/\tau_c$.
\end{gotcha}

\begin{note}
\textbf{Roadmap (v2).} The docstring enumerates the natural extensions, each of
which simply enlarges the auxiliary state space
(\file{pipeline/colored\_to\_markovian.py:51--64}): \emph{underdamped oscillatory}
kernels $\ee^{-|\tau|/\tau_d}\cos(\omega\tau)$ as a 2-state OU process with an
imaginary eigenvalue; \emph{double-exponential} kernels as a sum of two
independent OU processes; and \emph{polynomially-modulated} kernels as a chain of
OU processes (a higher-order linear filter). The single-Lorentzian case in this
chapter is the lowest-cost member of that family.
\end{note}

\section{Summary}

Colored noise is the one input that white-noise field theory cannot integrate
analytically. The fix is structural rather than numerical: rewrite the model so
the memory lives in an extra dynamical field. This chapter showed the full chain
--- a single-Lorentzian autocovariance \eqref{eq:cm-secondcumulant}, the OU filter
\eqref{eq:cm-toptline}--\eqref{eq:cm-bottomline} whose stationary covariance is
exactly that Lorentzian \eqref{eq:cm-ouvariance}, the three \msrjd{} structural
pieces and two action edits that encode it, the symbolic kernel matcher
\code{detect\_lorentzian} (\S\ref{sec:cm-detect}), and the in-place builder
rewrite \code{markovianize\_spec} (\S\ref{sec:cm-transform}) wired into
\code{build()} before compilation. The transformation is invisible to everything
downstream: by the time the field-theory core in the next part of the manual sees
your model, the colored noise has become an auxiliary OU field driven by white
noise, and the fast analytic integrators apply again.
